\section{LLM-Based Stakeholder Representation and Persona Modeling}
\label{sec:rw-persona}

Beyond RE automation, recent work has investigated how LLMs can represent users or stakeholders through synthetic agents and personas. This problem is particularly relevant to requirements elicitation because direct access to a sufficiently broad and diverse stakeholder population may be expensive, difficult, or unavailable.

Elicitron directly connects synthetic user simulation with requirements elicitation. It generates multiple LLM-based user agents, places them in simulated product-use scenarios, and subsequently interviews them to uncover explicit and latent user needs~\cite{ataei2025elicitron}. The study also compares alternative agent-generation methods and shows that the way simulated users are constructed affects both the diversity of generated agents and the needs discovered during elicitation.

Social Agents investigates persona-conditioned simulation at a broader population level. Its framework creates panels of agents conditioned on demographic and psychographic attributes, including factors such as age, location, interests, values, and lifestyle. Each persona agent generates an independent judgment, after which the individual outputs are aggregated according to a wisdom-of-crowds mechanism~\cite{bhattacharyya2026social}. The results distinguish structured persona diversity from stochastic repetition of an unconditioned LLM.

Beyond descriptive diversity, recent work increasingly examines whether simulated users exhibit realistic behavior. Persona Policies (PPol) addresses the tendency of default LLM user simulators to behave cooperatively and homogeneously. It uses evolutionary program search to generate task-preserving persona policies optimized for both human-likeness and behavioral coverage~\cite{chopra2026beyond}. The resulting personas vary along interactional properties such as communication style, information disclosure, clarification behavior, and reactions to errors.

PersonaGym focuses on another property of persona agents: whether an LLM consistently behaves according to its assigned persona. It introduces a dynamic evaluation framework and the PersonaScore metric for measuring persona adherence across persona-relevant environments~\cite{samuel2025personagym}. Its experiments demonstrate that general model capability or model size does not necessarily translate into stronger persona adherence.

Population-level representativeness introduces an additional challenge. \emph{LLM Generated Persona is a Promise with a Catch} examines large-scale LLM-generated persona populations and shows that adding generated persona information may amplify systematic bias and homogeneity rather than improve correspondence with the target population~\cite{li2026persona}. The study therefore emphasizes the need for more principled generation, calibration, and validation of synthetic personas.

MAJ-EVAL also provides a structured mechanism for constructing stakeholder-aligned personas. Rather than assigning arbitrary role descriptions, the framework first extracts stakeholder groups and evaluative perspectives from domain literature, and then constructs personas containing demographic information, evaluative dimensions, domain specialties, psychological traits, and social relationships~\cite{chen2026majeval}. These personas are subsequently instantiated as agents for multi-stakeholder evaluation.

Taken together, this literature shows that stakeholder or persona modeling involves more than assigning an agent a role label. Synthetic stakeholder representations may differ in how they are generated, grounded, behaviorally expressed, and validated. Consequently, a system that contains multiple stakeholder agents does not necessarily provide evidence that those agents constitute diverse, realistic, or representative stakeholder perspectives. A detailed comparison of these studies is given in Table~\ref{tab:apx-persona-studies} of Appendix~\ref{apx:persona-studies}.
